\section{Conclusion}
\label{sec:conclusion}

We evaluated adaptive multimodal evidence selection as a probability-dependent decision inside a fixed retrieval pipeline. Prior-aware calibration reduces the dominant probability error and allows expected-F1 to reach 0.1759 evidence F1 while returning 3.50 items per claim, a gain of 0.0085 over fixed top-5. The gain is concentrated in claims whose complete released-gold set is reachable. Calibrated thresholding remains stronger overall at 0.1807, while size and redundancy penalties add no measurable value. The shared verifier does not establish non-inferiority, and probability perturbations reverse the adaptive gain. Evidence-selector studies should therefore report three boundaries together: what retrieval makes reachable, whether membership scores are calibrated to the operating prior, and whether a structured decision rule improves on a calibrated threshold. In this setting, calibration adds value, extra decoder penalties do not, and retrieval remains the dominant limitation.
